
Weight space learning --- predicting model properties, enabling model merging, supporting continual learning --- depends on encoders that extract meaningful geometric structure from neural network weights. Two neural networks that compute the same function can differ by a scaling symmetry transform or a sign-flip symmetry transform, and may therefore occupy substantially different positions in weight space. If encoders operating on raw weights must contend with this variation, they may devote representational capacity to symmetry-induced variation rather than to functional properties.

The dominant paradigm, permutation-equivariant encoders [Zhou et al., 2023; Navon et al., 2023], addresses the neuron-relabeling symmetry of MLPs but leaves scaling and sign-flip symmetries --- which theory shows are equally well-defined for ReLU MLPs \citep{Navon2023EquivariantArchitectures} --- empirically uncharacterized in real model zoos. It remains unknown how large these orbits are in practice, whether existing equivariant encoders are naturally invariant to them, and whether canonicalization as a preprocessing step is effective.

This paper provides the first empirical quantification of scaling and sign-flip orbit diameters as cosine distances in a real model zoo --- distinct from prior empirical work [Entezari et al., 2022; Ainsworth et al., 2022] which characterizes permutation symmetry but not scaling or sign-flip orbit geometry. We measure these diameters in the Schürholt MNIST model zoo [Schürholt et al., 2022], a benchmark of 2-layer MLPs with ground-truth property labels, and we probe whether the Neural Functional Transformer (NFT) \citep{Zhou2023NeuralFunctionalNetworks} is naturally invariant to these orbits.

The results reveal a more nuanced picture than the theoretical argument suggests. Both symmetry types produce geometrically large orbits (scaling mean cosine distance 0.3232, sign-flip 1.075, 100\% above threshold), confirming the motivation for canonicalization. NFT's response is asymmetric: it is measurably non-invariant to scaling (gap=+0.024, CI entirely above zero) but is measurably invariant to sign-flip by construction (gap=-0.0007, CI entirely below zero). A structural audit reveals that the majority-sign sign-flip canonicalization algorithm is non-unique for 85.6\% of zoo models due to even d\_in=784, a combinatorial finding not previously reported in this literature. Property prediction experiments at N=500 are severely underpowered (CI width $\approx 0.6)$; all condition differences are statistically indistinguishable from zero.

We make four contributions:

\begin{enumerate}
\item \textbf{Orbit diameter measurement.} First empirical quantification of scaling and sign-flip orbit diameters as cosine distances in a real MLP zoo: scaling mean 0.3232 (CI=[0.3226, 0.3238]), sign-flip mean 1.075 (CI=[1.0705, 1.0789]), both 100\% above threshold across 2,500 oracle-constructed orbit pairs.
\end{enumerate}

\begin{enumerate}
\item \textbf{NFT invariance probe.} First quantification of NFT's non-invariance to scaling symmetry (gap=+0.024, CI=[0.0232, 0.0245]) and its significant approximate invariance to sign-flip symmetry (gap=-0.0007, CI=[-0.0008, -0.0005]) --- an asymmetric result not anticipated by design.
\end{enumerate}

\begin{enumerate}
\item \textbf{Geometric concentration.} Scaling canonicalization increases PCA explained variance ratio from 0.055 to 0.086 at k=20, confirming geometric redundancy removal. Linear probe $R^2$ and NFT Spearman $\rho$ do not improve at N=500.
\end{enumerate}

\begin{enumerate}
\item \textbf{Structural limitation characterization.} Majority-sign canonicalization produces a unique canonical form for only 14.4\% of Schürholt zoo models (d\_in=784, even), with failure rate accurately predicted by binomial combinatorics (predicted 83\%, observed 85.6\%).
\end{enumerate}

The remainder is organized as follows. Section 2 reviews related work. Section 3 describes the methodology. Section 4 presents the experimental setup. Section 5 reports results. Section 6 discusses implications and limitations. Section 7 concludes.

